% Generated by gen/generate.py; do not edit.
\begin{longtable}{@{}p{0.13\textwidth}p{0.07\textwidth}p{0.74\textwidth}@{}}
\toprule Statement & Lemma & Text \\ \midrule \endhead
\bottomrule \endfoot
Title & [1] & On the Meaning of Life \\[2pt]
Definition 1 & [2] & A life $L$ is \emph{meaningful to degree $d$} if $d$ is the extent to which $L$ consists in active engagement with projects of objective worth \citep{Wolf2010}. \\[2pt]
Assumption 2 & [3] & There exist projects of objective worth: projects whose value does not depend on the attitudes of those engaged in them \citep{Metz2013}. \\[2pt]
Assumption 3 & [4] & For every agent capable of intentional action, there is at least one project of objective worth in which she is able to engage. \\[2pt]
Theorem 4 & [5] & For every such agent, a life of positive meaning is available. \\[2pt]
Proof & [6] & By Assumption 3, there is a project $p$ of objective worth in which the agent can engage. \\[2pt]
 & [7] & A life that includes engagement with $p$ consists, to some positive extent, in active engagement with a project of objective worth. \\[2pt]
 & [8] & By Definition 1, it is meaningful to a positive degree. \\[2pt]
Remark 5 & [9] & The theorem answers the question of what the meaning of life is, and of whether meaning is available to everyone. \\[2pt]
 & [10] & It does not answer the question of whether any particular life is meaningful, which depends on what is done with it. \\[2pt]
 & [11] & Nor does it answer the sceptical challenge that, \emph{sub specie aeternitatis}, nothing matters \citep{Nagel1971}: Assumption 2 denies the premise of that challenge rather than refuting it. \\[2pt]
 & [12] & A defence of Assumption 2, and a discussion of the absurd, are given in the commentary (EXI-01-b). \\[2pt]
\end{longtable}
